%%\documentstyle[twocolumn,jsaiac]{jarticle}
%%\documentstyle[twocolumn,jsaiac]{j-article}
\documentclass[twocolumn]{jarticle}

\usepackage{jsaiac}

\usepackage{color}
\usepackage{url}
\usepackage{booktabs}
\usepackage{framed}

%%
\title{
\jtitle{実験指示から実行可能手順を生成するためのデータセット}
\etitle{A Dataset for Generating Executable Experimental Procedures from Experimental Instructions}
}
%%英文は以下を使用
%\title{Style file for manuscripts of JSAI 20XX}

\jaddress{山田 涼太，Science Aid株式会社，東京都中央区日本橋本町2丁目3番11号，ryota.yamada@science-aid.com}

\author{%
\jname{加藤 祥太\first}
\ename{Shota Kato}
\and
\jname{西田 理彦\second}
\ename{Toshihiko Nishida}
\and
\jname{酒井 雄介\third}
\ename{Yusuke Sakai}
\and
\jname{尾崎 遼\fourth}
\ename{Haruka Ozaki}
\and
\jname{杉山 亜矢斗\second}
\ename{Ayato Sugiyama}
\and
\jname{山田 涼太\fifth}
\ename{Ryota Yamada}
%Given-name Surname\third{}%%英文は左を使用
}

\affiliate{
\jname{\first{}京都大学}
\ename{Kyoto University}
\and
\jname{\second{}ラボラトリーオートメーション協会}
\ename{Laboratory Automation Suppliers' Association}
\and
\jname{\third{}東京大学}
\ename{The University of Tokyo}
\and
\jname{\fourth{}理化学研究所}
\ename{RIKEN}
\and
\jname{\fifth{}Science Aid株式会社}
\ename{Science Aid}
}

%%
%\Vol{28}        %% <-- 28th（変更しないでください）
%\session{0A0-00}%% <-- 講演ID（必須)

\begin{abstract}
実験指示から実行可能な実験手順を生成する能力を評価するため，ベンチマークデータセットLA-Bench 2025を構築し，LLM-as-a-judgeに基づく自動評価の設計とその特性を分析した．評価は共通採点基準（5点）と問題別の個別採点基準（5点）の合計10点満点とし，設計思想の異なる3種類の評価プロンプトを用いた．正解手順を対象に3種類のLLM評価器で各条件5回，計450回の採点を行った結果，プロンプト間で最大1.13点，評価器間で最大0.81点の得点差が生じた．この結果は，実験手順生成タスクの自動評価において，評価プロンプトの設計と評価器の選択が結果に直接影響することを示しており，単一の評価設計に依存しない評価手法の必要性を示唆する．
\end{abstract}

%\setcounter{page}{1}
\def\Style{``jsaiac.sty''}
\def\BibTeX{{\rm B\kern-.05em{\sc i\kern-.025em b}\kern-.08em%
 T\kern-.1667em\lower.7ex\hbox{E}\kern-.125emX}}
\def\JBibTeX{\leavevmode\lower .6ex\hbox{J}\kern-0.15em\BibTeX}
\def\LaTeXe{\LaTeX\kern.15em2$_{\textstyle\varepsilon}$}

\begin{document}
\maketitle

\section{はじめに}
自動化装置と機械学習を組み合わせ，実験の計画から実行・解析までを閉ループで回す自律実験は，材料・化学・生命科学の各分野で研究が進んでいる~\cite{hase2019}．この文脈では，「どの実験を次に実行するか」という探索・最適化戦略の性能を評価するベンチマークの整理が進んでいる．例えば，自己駆動ラボ（self-driving laboratory; SDL）のベンチマークに関するレビューでは，探索戦略をランダム探索のような参照手法と比較し，Acceleration factor（AF）および Enhancement factor（EF）によって定量化する枠組みが示されている~\cite{benchmarking_sdl_2026}．
一方で，自律実験の実装においては，探索よりも前段の工程として，目的や条件の記述といった実験指示を，自動化装置が実行可能な粒度の実験手順へ変換する工程が存在する．この工程では，物品制約，操作順序，数量整合，安全上の注意，参考手順との整合性など，多面的な制約を同時に満たす必要がある．本研究では，この「実験指示から実行可能手順を生成する能力の評価」を対象とする．

実験手順生成を対象とした評価は，人手評価が高コストであることに加え，合理的な別解が存在するため，単純な文字列一致や構造一致では妥当な評価が困難である．
Wet labプロトコルを対象とした情報抽出研究として，WLPデータセット~\cite{kulkarni2018wlp}やWNUT-2020共有タスク~\cite{tabassum2020wnut}，材料合成手順を構造化するPcMSP~\cite{yang2022pcmsp}がある．しかし，これらは既存手順の抽出・構造化を主目的とするものであり，「実験指示から新たに実行可能手順を生成する能力」の評価とは異なる．
近年は大規模言語モデル（large language model; LLM）を評価器として用いるLLM-as-a-judgeが提案されているが~\cite{zheng2023mtbench,liu2023geval}，実験手順のように安全性と実行可能性が重要な領域において，評価観点と評価プロンプトを明示的に設計し，評価対象のレイヤーを整理したベンチマークは存在しない．

本研究では，実験指示から実行可能な実験手順を生成する能力を評価するためのデータセット\textbf{LA-Bench 2025}を構築した．あわせて，実験実行可能性を構成する評価観点を明示的に定義し，各観点に対応した評価プロンプトを設計した．さらに，複数のLLMを評価器として用いた場合の評価結果の差異を分析し，評価器に依存した結果変動の可能性を示すことで，実験手順生成タスクにおける評価設計の重要性を明らかにする．

\section{方法}
\subsection{タスク定義とデータ構造}
LA-Bench 2025は，実験目的や条件を記述した自然言語の「実験指示」から，自動化装置で実行可能な粒度の「実験手順」を生成するタスクである．各問題は構造化された入力情報を持ち，それに対応する番号付き手順列を出力とする．入出力の構成を表\ref{tab:io_format}に示す．
入力の各要素は評価に直接用いられる．例えば，必ず使用する物品の一覧は手順中に適切に使用されていることが求められ，元プロトコルの手順は論理構造の整合性を評価する基準となる．また，期待される最終状態は，生成手順が実験として完遂可能であるかを判断するための参照情報として機能する．
出力手順は単なる構造化ではなく，暗黙知の補完，数値パラメータ（体積・濃度・温度・時間等）の明示，操作対象や装置の特定を含むことを求める．データ，評価コード，および詳細仕様は公開リポジトリに示す\footnote{\url{https://github.com/lasa-or-jp/la-bench}}．


\begin{table}[t]
    \centering
    \caption{LA-Bench 2025における入出力情報}
    \label{tab:io_format}
    \begin{tabular}{lp{6.5cm}}
    \toprule
    区分 & 内容 \\
    \midrule
    入力 & 1. 実験指示（実験目的と要求事項），\\
    & 2. 必ず使用する物品の一覧， \\
    & 3. 元プロトコルの手順（参考となる実験手順），\\
    & 4. 期待される最終状態（最終的に得られるサンプルや重要な中間生成物とその状態），\\
    & 5. 参考文献 \\
    \midrule
    出力 & 番号付きの実験手順．各手順は具体的な数量・条件・操作対象を明示し，学部レベルの訓練を受けた者が誤解なく実行可能な粒度とする．手順数は50以下，各手順は10文以下とする．\\
    \bottomrule
    \end{tabular}
\end{table}

LA-Bench 2025は例題セット5問，開発用テストセット10問，最終評価用テストセット9問からなる．各問題は，物品制約や論理構造（順序・分岐）を含む異なる操作特性を持つように設計した．本研究では開発用テストセットを主に用いる．

\subsection{評価設計}
実験手順生成の評価にあたり，本研究では「実行可能性」を入力情報の充足性，論理的一貫性，実験としての妥当性の三つの観点から捉える．単一指標では十分に評価できないため，観点ごとの得点を合算して総合点とする設計を採用した．

評価は，共通採点基準（5点満点）と問題ごとに定義される個別採点基準（5点満点）の合計10点満点で行う．

共通採点基準は，全問題に適用される一般的な実行可能性の評価である．具体的には，(1) 実験指示に含まれるパラメータが手順に正確に反映されているか，(2) 必須物品が適切に使用されているか，(3) 元プロトコルの論理構造（順序・分岐）が保持されているか，(4) 手順全体として期待される最終状態を論理的に達成可能か，(5) 明示されていない暗黙部分が適切に補完されているか，を評価する．
一方，計算誤り，論理矛盾，不自然な記述が認められる場合には減点する．減点は各項目につき1回のみ適用する．また，入力手順の単純転記や過度な安全側への条件拡張など，実質的な思考を伴わない出力と判断される場合には，共通採点基準の合計点を最大2点に制限する．この制限は「過度の安全性」による評価の歪みを抑制するためのものである．

個別採点基準は，各問題の出題意図に基づいて設定される．安全性，効率性，精度向上，コスト削減などの観点から具体的な達成条件を定義し，合計5点となるよう配点する．

\subsection{評価プロンプトと自動評価手順}
自動評価にはLLM-as-a-judgeを用いる．評価は，共通採点基準および個別採点基準に基づき，各回答に対して共通得点（0--5点）と個別得点（0--5点）を算出し，その合計を最終得点（0--10点）とする．また，共通基準の判定理由，個別基準に該当した根拠，および特記事項を併せて出力させる．

評価プロンプトは，システムプロンプトと，入力・出力・個別基準を埋め込む共通テンプレートから構成される．本研究では，評価器設計に起因する偏りを低減するため，設計思想の異なる3種類のシステムプロンプト（厳格なチェックリスト型，自律推論型，対照デュアルパス型）を独立に用い，それぞれで採点を行った．

厳格なチェックリスト型プロンプトでは，各共通採点基準を二値（OK/NG）で独立に判定し，不確実な場合は無加点とする保守的方針を採用した．判定根拠は出力手順からのみ引用させ，共通基準の判定理由を5行の固定フォーマットで出力させることで，得点と根拠の対応を明示的に拘束する．また，入力手順の単純転記や過度な安全側拡張が認められる場合には，共通採点基準の合計点を最大2点に制限する「過度の安全性cap」を適用する．

厳格なチェックリスト型プロンプトのうち，共通基準，過度の安全性cap，および出力形式に関する部分を図\ref{fig:strict_prompt}に示す．
\begin{figure}[!t]
  \centering
  \begin{minipage}{0.8\columnwidth}
  \begin{framed}
      \begin{verbatim}
# 共通採点基準（最大5点）
1. 実験指示のパラメータを正確に反映
2. 必須物品をすべて具体的に使用
3. 元プロトコルの順序・分岐を保持
4. 最終状態を論理的に達成可能
5. 暗黙部分を適切に補完

# 過度の安全性
- 入力手順の単純転記
- 出典の無批判な転記
- 不必要な安全側拡張
→ 共通得点を最大2点に制限

# 出力形式
{
  "general_score": 0-5,
  "specific_score": 0-5,
  "final_score": 0-10,
  "general_reason": "...",
  "specific_matches": [...],
  "notes": "..."
}
\end{verbatim}
  \vspace{-\baselineskip}
  \end{framed}
  \end{minipage}
  \caption{厳格なチェックリスト型評価プロンプトの抜粋．共通採点基準，過度の安全性cap，および出力形式を示す．}
  \label{fig:strict_prompt}
\end{figure}

自律推論型プロンプトでは，評価手順を明示し，自律的に推論を進める設計とした．対照デュアルパス型プロンプトでは，失敗理由探索と加点根拠探索を分離し，両者の衝突解決規則を明示することで，曖昧な判定を抑制した．各プロンプトの設計思想の比較は表\ref{tab:judge_prompt_comparison}に示す．
\begin{table*}[t]
    \centering
    \caption{3種類の評価プロンプトの設計思想の比較}
    \label{tab:judge_prompt_comparison}
    \begin{tabular}{llp{8cm}}
        \toprule
        プロンプト種別 & 設計の中心思想 & 主な特徴 \\ \midrule
        厳格なチェックリスト型 & 二値判定による保守的評価 &
        各共通採点基準を独立に判定し，不確実な場合は無加点とする．
        過度の安全性capを明示的に適用する． \\

        自律推論型（エージェント型） & 段階的推論による体系的評価 &
        入力抽出，出力解析，基準適用，集計の判断手順を明示して評価を行う．
        出力形式を制約し，冗長な出力を禁止する．\\

        対照デュアルパス型 & 失敗優先の対照評価 &
        失敗理由探索と加点根拠探索を分離する．
        衝突解決規則を明示し，曖昧な場合は無加点とする．\\ \bottomrule
    \end{tabular}
\end{table*}

確率的ばらつきを低減するため，同一入力に対して複数回の採点を行い，平均値を最終得点とした．さらに，複数のLLMを評価器として用い，得点差を比較することで評価器依存性を分析した．

評価プロンプト全文および実装の詳細は公開リポジトリを参照されたい．

\section{結果}
本節では，評価システムの特性を分析するため，開発用テストセット（10問）の正解手順（gold standard）を評価対象とした実験結果を報告する．評価器としてGPT-5-nano，GPT-5-mini，GPT-5の3種類のLLMを用い，3種類の評価プロンプトのそれぞれで各問題を5回ずつ採点した（temperature $= 1.0$，構造化出力を使用）．合計 $3 \times 3 \times 5 \times 10 = 450$ 回の評価を実施し，条件ごとの平均値を分析に用いた．評価器とプロンプトの組み合わせごとの結果を表\ref{tab:results_matrix}に示す．

\begin{table}[t]
    \centering
    \caption{評価器とプロンプトの組み合わせごとの最終得点平均}
    \label{tab:results_matrix}
    \begin{tabular}{lccc}
        \toprule
        評価器 & チェック & 自律推論 & デュアル \\
        & リスト型 & 型 & パス型 \\
        \midrule
        GPT-5-nano & 8.10 & 7.30 & 7.48 \\
        GPT-5-mini & 7.54 & 7.28 & 7.18 \\
        GPT-5      & 8.00 & 8.48 & 7.96 \\
        \midrule
        プロンプト平均 & 7.88 & 7.69 & 7.54 \\
        \bottomrule
    \end{tabular}
\end{table}

\subsection{評価スコアの分布}
全条件を通じた最終得点の平均は $\mu = 7.70$ 点，標準偏差は $\sigma = 2.30$ であった．問題ごとの平均得点（全条件にわたる平均）の範囲は $[3.42, 9.33]$ であり，満点（10点）に近い問題は限定的であった．

共通採点基準の平均は $3.49$ 点（5点満点），個別採点基準の平均は $4.22$ 点（5点満点）であった．個別基準が共通基準より高いことから，正解手順は問題固有の要件を概ね満たしている一方，共通基準が求める入力忠実性や論理的整合性の面では評価器により減点が生じやすいことが示された．

\subsection{評価プロンプト間の差異}
3種類の評価プロンプトにより得られた平均得点は，チェックリスト型で $\mu = 7.88$，自律推論型で $\mu = 7.69$，対照デュアルパス型で $\mu = 7.54$ であった（表\ref{tab:results_matrix}）．同一問題に対するプロンプト間の得点差は最大 $1.13$ 点であり，プロンプト設計が最終得点に影響することが確認された．

チェックリスト型が最も高い平均得点を示したのは，正解手順が各基準を明示的に満たしているため，二値判定において加点されやすいことが要因と考えられる．一方，対照デュアルパス型は失敗理由を優先的に探索する設計であるため，他のプロンプトでは加点される記述に対しても減点根拠を見出しやすく，最も低い得点となった．

これらの差異は，評価プロンプトの設計思想が評価結果に直接影響することを示している．

\subsection{評価器依存性}
評価器として異なるLLMを用いた場合，モデル間の全体平均得点差は $0.81$ 点であった（GPT-5: $8.15$，GPT-5-nano: $7.63$，GPT-5-mini: $7.33$）．プロンプトごとのモデル間得点差は，チェックリスト型で $0.56$ 点，自律推論型で $1.20$ 点，対照デュアルパス型で $0.78$ 点であり，自律推論型において最も大きかった．

GPT-5が最も高い得点を与える傾向は，より高い推論能力により正解手順の意図を適切に解釈し，加点根拠を見出しやすいことを反映している可能性がある．この結果は，単一モデルによる評価結果を絶対的な基準とすることは適切ではなく，複数評価器の併用が望ましいことを示唆する．

\section{考察}

本研究は，実験手順生成の評価において，評価観点の明示と評価プロンプト設計が重要であることを示した．評価結果はプロンプト設計および評価器の選択に依存し得るため，評価設計そのものが研究対象となる．

特に，共通採点基準における「過度の安全性cap」は，生成モデルが不必要に条件を安全側へ拡張する傾向を抑制するために導入した．この設計により，単なる条件の冗長化ではなく，入力忠実性と実行可能性のバランスを評価できるようになった．

一方で，LLM-as-a-judgeは自然言語的整合性を評価する能力に優れるが，数量整合や形式的制約の検証においては限界を持つ．今後は，形式的検証手法やルールベースの整合性チェックとLLM評価を統合する枠組みが有望である．

また，評価指標が最適化対象となると，モデルが指標に特化した挙動を学習する可能性がある．そのため，評価設計は固定的なものではなく，データセットの進化とともに更新されるべきである．

\section{おわりに}
本稿では，実験指示から実行可能な実験手順を生成するタスクに対して，LA-Bench 2025のデータ設計および評価設計を整理し，LLM-as-a-judgeに基づく採点プロンプトの構成を明示した．また，評価プロンプトおよび評価器の選択により得点が変動し得ることを示し，実験手順生成タスクにおける評価設計の重要性を述べた．

本データセットはコンペティションでも使用され，複数のグループから提出された生成手順に対して，自動評価および一部の提出に対する専門家による人手評価が実施された．しかし，本稿では主としてデータと自動評価設計に焦点を当てたため，提出物全体の傾向分析や，人手評価結果と自動評価結果の差異分析，別解の扱いを含む評価の妥当性検証は今後の課題である．

\section*{謝辞}
LA-Bench 2025に参加し，実験手順生成手法を提出いただいた各チームの皆様に感謝する．
また，本取り組みは一般社団法人ラボラトリーオートメーション協会（LASA）および人工知能学会のコンペティション開催支援制度の支援を受けた．さらに，スポンサーとして株式会社ディープコア，合同会社embrio，エピストラ株式会社，HERO Impact Capital，株式会社Inner Resource，株式会社システム計画研究所，株式会社リンカーズOI研究所，Polaris.AI株式会社，プライマルキャピタルの各団体にご支援をいただいた．ここに記して感謝する．

\bibliographystyle{jsai}
\bibliography{reference}

\end{document}
